\chapter{Vim Reference Card}\label{app:vim}

\section{Modes}\label{app:vim-modes}

\begin{keyslong}[0.24]
	\key{<Esc>}   & Return to Normal mode from anywhere        \\
	\key{i} \key{a} \key{I} \key{A} \key{o} \key{O} & Enter Insert mode \\
	\key{v} \key{V} \key{<C-v>} & Charwise / linewise / blockwise Visual \\
	\key{:}       & Command-line mode                          \\
	\key{R}       & Replace mode                               \\
	\key{<C-\textbackslash>} \key{<C-n>} & Leave Terminal mode \\
\end{keyslong}

\section{Motions}\label{app:vim-motions}

\begin{keyslong}[0.24]
	\key{h} \key{j} \key{k} \key{l}   & Left, down, up, right                  \\
	\key{w} \key{b} \key{e}           & Word forward, back, end                \\
	\key{W} \key{B} \key{E}           & Same, by whitespace-delimited WORD     \\
	\key{0} \key{\textasciicircum{}} \key{\$}         & Line start, first non-blank, line end  \\
	\key{f} \key{x} / \key{t} \key{x} & To / till the next \texttt{x}; \key{;} and \key{,} repeat \\
	\key{g} \key{g} / \key{G}         & Top / bottom of the file               \\
	\key{42} \key{G}                  & Line 42                                \\
	\key{\{} \key{\}}                 & Previous / next paragraph              \\
	\key{\%}                          & Matching bracket                       \\
	\key{H} \key{M} \key{L}           & Top / middle / bottom of the screen    \\
	\key{<C-d>} \key{<C-u>}           & Half page down / up                    \\
	\key{z} \key{z}                   & Centre the cursor line                 \\
	\key{<C-o>} \key{<C-i>}           & Back / forward in the jump list        \\
	\key{m} \key{a} then \key{`} \key{a} & Set and return to mark \texttt{a}   \\
\end{keyslong}

\section{Operators and Text Objects}\label{app:vim-operators}

\begin{keyslong}[0.24]
	\key{d} \key{c} \key{y}           & Delete, change, yank                   \\
	\key{>} \key{<} \key{=}           & Indent, dedent, re-indent              \\
	\key{g} \key{u} \key{g} \key{U} \key{g} \key{\textasciitilde} & Lower, upper, toggle case \\
	\key{g} \key{q}                   & Reformat to \cmd{'textwidth'}          \\
	\key{d} \key{d} \key{y} \key{y} \key{c} \key{c} & Operate on the whole line \\
	\key{D} \key{C} \key{Y}           & Operate to the end of the line         \\
	\key{x} / \key{X}                 & Delete the character under / before the cursor \\
	\key{r} \key{x}                   & Replace one character with \texttt{x}  \\
	\key{J}                           & Join with the next line                \\
	\key{i} \key{w} / \key{a} \key{w} & Inner / a word                         \\
	\key{i} \key{p} / \key{a} \key{p} & Inner / a paragraph                    \\
	\key{i} \key{"} \key{i} \key{(} \key{i} \key{\{} \key{i} \key{[} \key{i} \key{t} & Inside quotes, parens, braces, brackets, tags \\
	\key{a} \key{"} \key{a} \key{(} \dots & The same, including the delimiters \\
	\key{.}                           & Repeat the last change                 \\
	\key{u} / \key{<C-r>}             & Undo / redo                            \\
\end{keyslong}

\section{Registers, Macros and Search}\label{app:vim-registers}

\begin{keyslong}[0.24]
	\key{"} \key{a} \key{y} \key{y}   & Yank a line into register \texttt{a}   \\
	\key{"} \key{a} \key{p}           & Paste register \texttt{a}              \\
	\key{"} \key{0} \key{p}           & Paste the last \emph{yank}, not the last delete \\
	\key{"} \key{+} \key{y} / \key{"} \key{+} \key{p} & System clipboard      \\
	\key{"} \key{\_} \key{d}          & Delete without touching any register   \\
	\cmd{:registers}                  & Show every register                    \\
	\key{q} \key{a} \dots \key{q}     & Record a macro into register \texttt{a} \\
	\key{@} \key{a} / \key{@} \key{@} & Play it / replay the last one          \\
	\key{/} \texttt{pat} \key{<CR>}   & Search forward; \key{n} and \key{N} repeat \\
	\key{*} / \key{\#}                & Search for the word under the cursor   \\
	\cmd{:noh}                        & Clear search highlighting              \\
	\cmd{:\%s/old/new/gc}             & Substitute in the file, confirming each \\
	\cmd{:g/pat/normal @q}            & Run macro \texttt{q} on matching lines \\
\end{keyslong}

\section{Files, Buffers and Windows}\label{app:vim-buffers}

\begin{keyslong}[0.24]
	\cmd{:w} \cmd{:q} \cmd{:wq} \cmd{:q!} & Write, quit, both, force-quit      \\
	\cmd{:e file} / \cmd{:ls} / \cmd{:b N} & Open, list, switch buffer         \\
	\key{<C-\textasciicircum{}>}                      & Toggle with the alternate buffer       \\
	\key{<C-w>} \key{s} / \key{<C-w>} \key{v} & Split horizontally / vertically \\
	\key{<C-w>} \key{h} \key{j} \key{k} \key{l} & Move between windows         \\
	\key{<C-w>} \key{c} / \key{<C-w>} \key{o} & Close this / all other windows \\
	\cmd{:tabnew} / \key{g} \key{t}   & New tab / next tab                     \\
	\cmd{:copen} / \cmd{:cn} / \cmd{:cp} & Quickfix list: open, next, previous \\
	\cmd{:\%!cmd}                     & Filter the buffer through a shell command \\
	\cmd{:help <topic>}               & The manual                             \\
\end{keyslong}

\section{Neovim}\label{app:vim-neovim}

\begin{keyslong}[0.30]
	\cmd{:checkhealth}                & Diagnose the whole installation        \\
	\cmd{:Lazy}                       & Plugin manager dashboard               \\
	\cmd{:Mason}                      & Install language servers and formatters \\
	\cmd{:LspInfo}                    & Which servers serve this buffer        \\
	\cmd{:Inspect} / \cmd{:InspectTree} & Highlight groups / Treesitter tree   \\
	\cmd{:terminal}                   & Shell in a buffer                      \\
	\cmd{:lua} \dots{} / \cmd{:=expr} & Execute / print Lua                    \\
	\cmd{:verbose set opt?}           & Which file set this option             \\
	\cmd{:source \%}                  & Re-run the Lua file being edited       \\
	\key{K} \key{g} \key{d} \key{g} \key{r} & LSP hover, definition, references \\
	\key{]} \key{d} / \key{[} \key{d} & Next / previous diagnostic             \\
\end{keyslong}

\chapter{Shell Reference Card}\label{app:sh}

\section{Files and Navigation}\label{app:sh-files}

\begin{keyslong}[0.30]
	\cmd{pwd}, \cmd{cd d}, \cmd{cd -}   & Where am I; go there; go back             \\
	\cmd{ls -lh}, \cmd{ls -la}          & Long listing; include hidden files        \\
	\cmd{tree -L 2}                     & Two levels of the tree                    \\
	\cmd{mkdir -p a/b/c}                & Make a whole chain of directories         \\
	\cmd{cp -r src dst}, \cmd{mv a b}   & Copy a tree; move or rename               \\
	\cmd{rm -r d}, \cmd{rm -i f}        & Delete a tree; ask first                  \\
	\cmd{ln -s target name}             & Symbolic link                             \\
	\cmd{chmod +x f}, \cmd{chmod 644 f} & Make executable; set exact permissions    \\
	\cmd{du -sh *}, \cmd{df -h}         & Size of each entry; free space            \\
	\cmd{stat f}, \cmd{file f}          & Metadata; what it actually is             \\
	\cmd{find . -name '*.c' -type f}    & Find by name --- quote the pattern        \\
	\cmd{find . -mtime -7 -size +1M}    & Recent and large                          \\
	\cmd{find . -name '*.o' -delete}    & Find and remove                           \\
	\cmd{find . -name '*.c' -exec cmd \{\} +} & Run a command on the results        \\
\end{keyslong}

\section{Text Processing}\label{app:sh-text}

\begin{keyslong}[0.30]
	\cmd{cat}, \cmd{head -n 3}, \cmd{tail -f} & Print all; first three; follow      \\
	\cmd{less f}                        & Page through it; Vim keys apply           \\
	\cmd{wc -l f}                       & Count lines                               \\
	\cmd{grep -rn pat dir}              & Search recursively, with line numbers     \\
	\cmd{grep -i}, \cmd{-v}, \cmd{-w}, \cmd{-c} & Ignore case; invert; whole word; count \\
	\cmd{grep -E 'a|b'}                 & Extended regular expressions              \\
	\cmd{grep -q pat f}                 & Silent --- for use as a test              \\
	\cmd{sed 's/old/new/g'}             & Substitute everywhere on each line        \\
	\cmd{sed -n '10,20p'}               & Print a range of lines                    \\
	\cmd{sed -i.bak 's/a/b/g' f}        & Edit in place, keeping a backup           \\
	\cmd{awk '\{print \$2\}'}           & The second field of every line            \\
	\cmd{awk -F: '\$3 > 999 \{print \$1\}'} & Field test with a custom separator    \\
	\cmd{sort -rn}, \cmd{sort -u}, \cmd{sort -h} & Numeric descending; unique; by size \\
	\cmd{uniq -c}                       & Count adjacent duplicates                 \\
	\cmd{cut -d: -f1}, \cmd{tr -d '\textbackslash r'} & Field one; strip carriage returns \\
	\cmd{diff -u a b}                   & Unified diff --- the format Git uses      \\
\end{keyslong}

\section{Redirection and Pipes}\label{app:sh-pipes}

\begin{keyslong}[0.30]
	\texttt{> f}, \texttt{>{}> f}         & Overwrite; append                         \\
	\texttt{< f}                        & Read standard input from a file           \\
	\texttt{2> f}, \texttt{2>\&1}, \texttt{\&> f} & Errors to a file; errors to wherever output goes; both \\
	\texttt{> /dev/null 2>\&1}          & Discard everything                        \\
	\texttt{a | b}                      & Connect output to input                   \\
	\cmd{cmd | tee log}                 & Save and pass through                     \\
	\cmd{... | xargs -0 rm}             & Turn input into arguments, NUL-separated  \\
	\texttt{\$(cmd)}                    & Command substitution --- always quote it  \\
	\texttt{<(cmd)}                     & Process substitution: output as a filename \\
	\texttt{<{}<'EOF' ... EOF}            & Here-document, nothing expanded           \\
	\texttt{a \&\& b}, \texttt{a || b}  & Run on success; run on failure            \\
	\texttt{\$?}                        & Exit status of the last command           \\
\end{keyslong}

\section{Processes and Jobs}\label{app:sh-proc}

\begin{keyslong}[0.30]
	\cmd{ps aux}, \cmd{pgrep n}, \cmd{pstree} & All processes; by name; as a tree   \\
	\cmd{top}, \cmd{htop}               & Live resource use                         \\
	\cmd{lsof -i :8080}                 & What is holding this port                 \\
	\cmd{cmd \&}, \key{<C-z>}, \cmd{bg}, \cmd{fg} & Background; suspend; resume there; resume here \\
	\cmd{jobs -l}                       & This shell's jobs, with PIDs              \\
	\cmd{kill PID}, \cmd{kill -9 PID}   & Polite; last resort                       \\
	\cmd{pkill -f pattern}              & Kill by command line                      \\
	\cmd{nohup cmd \&}, \cmd{disown}    & Survive the shell exiting                 \\
	\cmd{export VAR=v}, \cmd{env}       & Set for children; show the environment    \\
	\cmd{VAR=v cmd}                     & Set for one command only                  \\
	\cmd{echo "\$PATH"}, \cmd{which -a} & Search path; every match in order         \\
	\cmd{source f}                      & Run in the current shell, not a child     \\
	\cmd{type x}                        & Builtin, alias, function, or file?        \\
	\cmd{time cmd}                      & Wall, user and system time                \\
\end{keyslong}

\section{Scripting Idioms}\label{app:sh-script}

\begin{keyslong}[0.30]
	\texttt{\#!/usr/bin/env bash}       & Portable shebang                          \\
	\cmd{set -euo pipefail}             & Fail on error, unset variable, or broken pipe \\
	\cmd{trap 'rm -rf "\$tmp"' EXIT}    & Clean up however the script ends          \\
	\texttt{"\$@"}                      & All arguments, each as its own word       \\
	\texttt{\$\#}, \cmd{shift}          & How many; discard the first               \\
	\texttt{\$\{v:-default\}}           & Fall back if unset or empty               \\
	\texttt{\$\{f\#\#*/\}}, \texttt{\$\{f\%/*\}} & Basename; dirname                \\
	\texttt{\$\{f\%.tex\}.pdf}          & Swap an extension                         \\
	\texttt{[[ -f f ]]}, \texttt{[[ -z "\$v" ]]} & File exists; string is empty     \\
	\texttt{[[ \$a == \$b ]]}, \texttt{[[ \$n -lt 5 ]]} & String and numeric comparison \\
	\cmd{for f in *.tex; do ...; done}  & Loop over files --- never over \cmd{ls}   \\
	\cmd{while IFS= read -r line}       & Read lines exactly as written             \\
	\cmd{local v=\$1}                   & Keep a function's variables to itself     \\
	\cmd{mktemp -d}                     & A temporary directory, safely named       \\
	\cmd{bash -x script.sh}             & Trace every command after expansion       \\
	\cmd{shellcheck script.sh}          & Static analysis                           \\
\end{keyslong}

\chapter{Make Reference Card}\label{app:mk}

\section{Rule Syntax}\label{app:mk-rules}

\begin{keyslong}[0.34]
	\texttt{target: prereqs}          & A rule; the recipe follows, tab-indented   \\
	\texttt{\textrightarrow recipe}   & Recipe lines begin with a literal TAB      \\
	\texttt{target: prereqs | order-only} & After \texttt{|}: built first, timestamp ignored \\
	\texttt{\%.o: \%.c}               & Pattern rule; \texttt{\%} is the stem      \\
	\texttt{\$(OBJ): \%.o: \%.c}      & Static pattern rule, restricted to a list  \\
	\texttt{.PHONY: all clean}        & Targets that are not files                 \\
	\texttt{.DEFAULT\_GOAL := app}    & What a bare \cmd{make} builds              \\
	\texttt{vpath \%.c src}           & Where to look for matching prerequisites   \\
	\texttt{-include \$(DEP)}         & Include a file, silently if it is missing  \\
	\texttt{@cmd}                     & Do not echo this recipe line               \\
	\texttt{-cmd}                     & Ignore a non-zero exit status              \\
	\texttt{cmd1 \&\& cmd2}           & One shell: each recipe line gets its own   \\
	\texttt{SHELL := /bin/bash}       & Use Bash instead of \file{/bin/sh}         \\
\end{keyslong}

\section{Variables}\label{app:mk-vars}

\begin{keyslong}[0.34]
	\texttt{VAR = value}       & Recursive: re-expanded at every use               \\
	\texttt{VAR := value}      & Simple: expanded once, here. Prefer this          \\
	\texttt{VAR ?= value}      & Only if not already set --- lets callers override \\
	\texttt{VAR += value}      & Append                                            \\
	\texttt{VAR != cmd}        & Assign a shell command's output                   \\
	\texttt{target: VAR += x}  & Target-specific value, inherited by prerequisites \\
	\texttt{override VAR += x} & Wins even over the command line                   \\
	\texttt{\$\$}              & A literal \texttt{\$}, passed on to the shell     \\
	\texttt{CC CFLAGS CPPFLAGS} & Compiler, its flags, preprocessor flags          \\
	\texttt{LDFLAGS LDLIBS}    & Where to look; what to link. \cmd{-lm} goes in \texttt{LDLIBS} \\
	\texttt{PREFIX DESTDIR}    & Install location; packager's staging root         \\
	\texttt{\$(MAKE)}          & Recursive invocation --- never write \cmd{make}   \\
\end{keyslong}

\section{Automatic Variables}\label{app:mk-auto}

\begin{keyslong}[0.20]
	\texttt{\$@}   & The target                                                     \\
	\texttt{\$<}   & The first prerequisite --- use for \cmd{-c}                    \\
	\texttt{\$\textasciicircum{}} & All prerequisites, deduplicated --- use for linking \\
	\texttt{\$+}   & All prerequisites, duplicates kept                             \\
	\texttt{\$?}   & Only those newer than the target                               \\
	\texttt{\$*}   & The stem matched by \texttt{\%}                                \\
	\texttt{\$(@D)}, \texttt{\$(@F)} & The target's directory and filename          \\
\end{keyslong}

\section{Built-in Functions}\label{app:mk-funcs}

\begin{keyslong}[0.38]
	\texttt{\$(wildcard *.c)}               & Expand a glob now                     \\
	\texttt{\$(patsubst \%.c,\%.o,\$(SRC))} & Pattern substitution                  \\
	\texttt{\$(SRC:.c=.o)}                  & The same, in shorthand                \\
	\texttt{\$(subst a,b,\$(T))}            & Plain text substitution               \\
	\texttt{\$(filter \%.c,\$(F))}          & Keep matching words                   \\
	\texttt{\$(filter-out main.o,\$(O))}    & Drop matching words                   \\
	\texttt{\$(sort \$(L))}                 & Sort and deduplicate                  \\
	\texttt{\$(dir \$(F))}, \texttt{\$(notdir \$(F))} & Directory, filename         \\
	\texttt{\$(basename \$(F))}, \texttt{\$(suffix \$(F))} & Without, only the extension \\
	\texttt{\$(addprefix build/,\$(O))}     & Prepend to every word                 \\
	\texttt{\$(shell date +\%F)}            & Run a shell command                   \\
	\texttt{\$(info \dots)}                 & Print while reading the file          \\
	\texttt{\$(error \dots)}                & Print and stop                        \\
\end{keyslong}

\section{Command-Line Options}\label{app:mk-flags}

\begin{keyslong}[0.30]
	\cmd{make -n}          & Dry run: print the recipes, run none                   \\
	\cmd{make -j\$(nproc)} & One job per core                                       \\
	\cmd{make -B}          & Rebuild everything regardless of timestamps            \\
	\cmd{make -k}          & Carry on after errors                                  \\
	\cmd{make -C dir}      & Change directory first                                 \\
	\cmd{make -f other.mk} & Use a different file                                   \\
	\cmd{make -r}          & Drop the built-in rule database                        \\
	\cmd{make -s}          & Do not echo recipes                                    \\
	\cmd{make --debug=b}   & Say why each target was or was not rebuilt             \\
	\cmd{make --trace}     & Annotate each recipe with its Makefile line            \\
	\cmd{make -p}          & Dump every rule and variable                           \\
	\cmd{make CC=clang}    & Override a variable for this build                     \\
\end{keyslong}

\section{Diagnosing a Makefile}\label{app:mk-debug}

\begin{keyslong}[0.34]
	\texttt{missing separator}        & Spaces where a tab belongs. \cmd{cat -A Makefile} \\
	\texttt{No rule to make target}   & Misspelt filename, or a missing generator rule \\
	\texttt{'x' is up to date}        & \texttt{x} should have been in \texttt{.PHONY} \\
	\texttt{Recursive variable \dots} & \texttt{=} where \texttt{:=} was meant      \\
	Fails only under \cmd{-j}         & A missing edge in the graph, not a make bug \\
	Rebuilds every time               & A directory used as a normal prerequisite; make it order-only \\
	Does not rebuild on a header edit & Add \cmd{-MMD -MP} and \cmd{-include \$(DEP)} \\
	\cmd{print-\%:} \texttt{@echo '\$* = [\$(\$*)]'} & A rule for printing any variable \\
\end{keyslong}

\chapter{GDB Reference Card}\label{app:gdb}

\section{Starting and Stopping}\label{app:gdb-start}

\begin{keyslong}[0.32]
	\cmd{gcc -g -O0 -o app *.c}   & Compile so the binary can be debugged at all   \\
	\cmd{gdb ./app}               & Load it; nothing runs yet                      \\
	\cmd{gdb -{}-args ./app -v f} & \dots{} with arguments                         \\
	\cmd{gdb ./app core}          & Post-mortem from a core dump                   \\
	\cmd{gdb -p 1234}             & Attach to a running process                    \\
	\cmd{gdb -q -batch -ex run -ex bt -{}-args ./app} & One-line crash triage      \\
	\cmd{start}                   & Run, stopping at the first line of \cmd{main}  \\
	\cmd{run} / \cmd{r}           & Run from the beginning                         \\
	\cmd{continue} / \cmd{c}      & Resume                                         \\
	\cmd{kill}, \cmd{quit}        & Stop the program; leave GDB                    \\
	\cmd{set args ...}            & Arguments for the next \cmd{run}               \\
	\key{<CR>}                    & Repeat the previous command                    \\
	\key{<C-c>}                   & Interrupt the running program                  \\
	\cmd{help break}, \cmd{apropos watch} & Help on one command; search all help   \\
\end{keyslong}

\section{Breakpoints and Watchpoints}\label{app:gdb-break}

\begin{keyslong}[0.32]
	\cmd{break main}                 & On a function                               \\
	\cmd{break parser.c:42}          & On a line                                   \\
	\cmd{break parse if i == 500}    & Conditional --- the one that saves the day  \\
	\cmd{tbreak main}                & Temporary: deleted after it fires           \\
	\cmd{ignore 2 499}               & Skip the next 499 hits. Faster than a condition \\
	\cmd{condition 2 n > 100}        & Add or change a condition                   \\
	\cmd{info breakpoints}           & List them, with hit counts                  \\
	\cmd{disable 2}, \cmd{enable 2}, \cmd{delete 2} & Manage one                   \\
	\cmd{watch total}                & Stop when a value changes                   \\
	\cmd{watch -l *ptr}              & Watch a location, not a name                \\
	\cmd{rwatch}, \cmd{awatch}       & Stop on read; on read or write              \\
	\cmd{catch throw}, \cmd{catch syscall write} & Exceptions; system calls        \\
	\cmd{commands} \dots{} \cmd{end} & Script to run at a breakpoint               \\
	\cmd{silent} / \cmd{printf} / \cmd{continue} & The usual trio inside one       \\
\end{keyslong}

\section{Stepping and Frames}\label{app:gdb-step}

\begin{keyslong}[0.32]
	\cmd{next} / \cmd{n}      & Next source line, stepping \emph{over} calls       \\
	\cmd{step} / \cmd{s}      & Next source line, stepping \emph{into} calls       \\
	\cmd{finish}              & Run until this function returns                    \\
	\cmd{until}               & Like \cmd{next}, but run a loop to completion      \\
	\cmd{until 60}, \cmd{advance f} & Run to a line, or to a function              \\
	\cmd{stepi}, \cmd{nexti}  & One machine instruction                            \\
	\cmd{return 0}            & Abandon the frame, forcing a return value          \\
	\cmd{jump 42}             & Continue from another line in this function        \\
	\cmd{backtrace} / \cmd{bt} & The call stack. The first command after a crash   \\
	\cmd{bt full}             & \dots{} with every frame's locals                  \\
	\cmd{frame 2}, \cmd{up}, \cmd{down} & Select a frame                           \\
	\cmd{info locals}, \cmd{info args} & The selected frame's variables            \\
	\cmd{list}, \cmd{list parse} & Show source                                     \\
	\cmd{disassemble /s}      & Instructions, interleaved with source              \\
\end{keyslong}

\section{Examining Data}\label{app:gdb-data}

\begin{keyslong}[0.32]
	\cmd{print n} / \cmd{p n}   & Evaluate an expression in the current frame      \\
	\cmd{print/x}, \cmd{/t}, \cmd{/c}, \cmd{/s} & Hex, binary, character, string   \\
	\cmd{print *p@n}            & \texttt{n} elements from a pointer. Essential in C \\
	\cmd{print *cfg}            & Dereference a struct pointer                     \\
	\cmd{print func(3)}         & Call a function in the live process              \\
	\cmd{x/16xb buf}            & 16 bytes of memory, in hex                       \\
	\cmd{x/s s}, \cmd{x/5i \$pc} & As a string; as instructions                    \\
	\cmd{ptype cfg}, \cmd{whatis cfg} & The type, expanded or named                \\
	\cmd{ptype/o struct node}   & \dots{} with offsets and padding                 \\
	\cmd{display i}             & Print it automatically at every stop             \\
	\cmd{set var n = 42}        & Assign --- note \cmd{var}, not bare \cmd{set}    \\
	\cmd{set \$p = head}        & A convenience variable                           \\
	\cmd{\$1}, \cmd{\$}         & Earlier results; the last one                    \\
	\cmd{info registers}        & Every register                                   \\
\end{keyslong}

\section{Threads, Signals and Comfort}\label{app:gdb-threads}

\begin{keyslong}[0.34]
	\cmd{info threads}              & List them                                    \\
	\cmd{thread apply all bt}       & Every thread's stack. The deadlock command   \\
	\cmd{thread 3}                  & Switch                                       \\
	\cmd{set scheduler-locking on}  & Step one thread, freeze the others           \\
	\cmd{handle SIGPIPE nostop noprint pass} & Stop GDB interrupting on a routine signal \\
	\cmd{info signals}              & The whole signal table                       \\
	\key{<C-x>} \key{a}             & Toggle the TUI                               \\
	\cmd{layout src}, \cmd{focus cmd} & Choose a layout; give the arrows back      \\
	\cmd{set print pretty on}       & Readable structs                             \\
	\cmd{set pagination off}        & No ``---Type <return> to continue---''       \\
	\cmd{define} \dots{} \cmd{end}  & A user-defined command                       \\
	\cmd{source debug.gdb}          & Load a file of commands                      \\
	\cmd{:packadd termdebug} then \cmd{:Termdebug ./app} & GDB inside Neovim       \\
\end{keyslong}

\section{Choosing the Right Tool}\label{app:gdb-triage}

\begin{keyslong}[0.36]
	Crash, sensible backtrace        & GDB: \cmd{bt}, \cmd{info args}, \cmd{up}    \\
	Crash, backtrace is \texttt{??}  & The stack is corrupt --- a sanitiser first  \\
	Use-after-free, buffer overrun   & \cmd{gcc -fsanitize=address}                \\
	Signed overflow, bad shift       & \cmd{gcc -fsanitize=undefined}              \\
	Data race                        & \cmd{gcc -fsanitize=thread}                 \\
	Hang or deadlock                 & \cmd{gdb -p} then \cmd{thread apply all bt} \\
	Wrong answer, no crash           & A breakpoint and \cmd{display}              \\
	Wrong value, unknown writer      & \cmd{watch}                                 \\
	Happens on iteration 500         & \cmd{break f if i == 500}                   \\
	Only fails at \cmd{-O2}          & Suspect undefined behaviour, not the compiler \\
	Crashed once, on a server        & \cmd{ulimit -c unlimited}, then \cmd{gdb ./app core} \\
\end{keyslong}

\chapter{Git Reference Card}\label{app:git}

\section{The Everyday Loop}\label{app:git-loop}

\begin{keyslong}[0.34]
	\cmd{git init}, \cmd{git clone <url>} & Start a repository; copy one           \\
	\cmd{git status -sb}                & What changed, in two columns             \\
	\cmd{git add file}, \cmd{git add -A} & Stage one file; stage everything        \\
	\cmd{git add -p}                    & Stage parts of files, interactively      \\
	\cmd{git rm -{}-cached file}        & Stop tracking, keep on disk              \\
	\cmd{git commit -v}                 & Commit, with the diff in the editor      \\
	\cmd{git commit -m "..."}           & Commit with a one-line message           \\
	\cmd{git commit -{}-amend}          & Redo the last commit                     \\
	\cmd{git diff}                      & Working tree vs index (unstaged)         \\
	\cmd{git diff -{}-staged}           & Index vs \texttt{HEAD} (staged)          \\
	\cmd{git diff HEAD}                 & Everything not committed                 \\
	\cmd{git diff -w}, \cmd{-{}-word-diff} & Ignore whitespace; diff by word       \\
\end{keyslong}

\section{Branching and Merging}\label{app:git-branch}

\begin{keyslong}[0.34]
	\cmd{git switch -c feature}          & Create a branch and move there           \\
	\cmd{git switch main}, \cmd{git switch -} & Move; move back                     \\
	\cmd{git branch -v}, \cmd{git branch -d f} & List with commits; delete          \\
	\cmd{git merge feature}              & Merge into the current branch            \\
	\cmd{git merge -{}-no-ff feature}    & Always record a merge commit             \\
	\cmd{git merge -{}-abort}            & Undo an in-progress merge                \\
	\cmd{git rebase main}                & Replay this branch on top of \cmd{main}  \\
	\cmd{git rebase -i main}             & \dots{} editing the series first         \\
	\cmd{git rebase -{}-continue} / \cmd{-{}-abort} & After a conflict; give up     \\
	\cmd{pick} / \cmd{reword} / \cmd{squash} / \cmd{fixup} / \cmd{drop} & The interactive-rebase verbs \\
	\cmd{git commit -{}-fixup=<sha>}     & Mark a fix, for \cmd{-{}-autosquash}     \\
	\cmd{git cherry-pick <sha>}          & Apply one commit here                    \\
	\cmd{git stash}, \cmd{git stash pop} & Set work aside; bring it back            \\
	\cmd{git worktree add ../hotfix main} & A second directory, another branch      \\
\end{keyslong}

\section{Remotes}\label{app:git-remote}

\begin{keyslong}[0.36]
	\cmd{git remote -v}                  & Which remotes, and their URLs            \\
	\cmd{git remote add upstream <url>}  & Add one --- the fork convention          \\
	\cmd{git fetch}, \cmd{git fetch -{}-prune} & Download; forget dead branches     \\
	\cmd{git log HEAD..origin/main}      & What they did that you do not have       \\
	\cmd{git pull -{}-rebase}            & Fetch, then replay your commits on top   \\
	\cmd{git push -u origin feature}     & Push and set the upstream                \\
	\cmd{git push -{}-force-with-lease}  & Overwrite, but only if nobody else pushed \\
	\cmd{git push origin -{}-delete f}   & Delete a remote branch                   \\
	\cmd{git tag -a v1.0 -m "..."}       & An annotated tag                         \\
	\cmd{git push -{}-follow-tags}       & Push commits with their tags             \\
	\cmd{git describe -{}-tags}          & Nearest tag, plus distance               \\
	\cmd{gh pr create -{}-fill}          & Open a pull request                      \\
	\cmd{gh pr checkout 42}              & Check out someone's PR to actually run it \\
	\cmd{gh pr checks 42}, \cmd{gh run watch} & Did CI pass; follow a run live      \\
\end{keyslong}

\section{Inspecting History}\label{app:git-history}

\begin{keyslong}[0.36]
	\cmd{git log -{}-oneline -{}-graph -{}-all -{}-decorate} & The one to alias     \\
	\cmd{git log -p -{}- file.tex}       & Full diffs, one file                     \\
	\cmd{git log -S'fopen'}              & Commits that added or removed a string   \\
	\cmd{git log -L 40,60:file.c}        & The history of specific lines            \\
	\cmd{git log main..feature}          & On one branch but not the other          \\
	\cmd{git show <sha>}                 & One commit in full                       \\
	\cmd{git blame -w -C file.c}         & Who last touched each line               \\
	\cmd{git bisect start / bad / good}  & Binary-search for the breaking commit    \\
	\cmd{git bisect run ./test.sh}       & \dots{} unattended                       \\
	\cmd{git cat-file -p HEAD}           & The raw commit object                    \\
	\cmd{git count-objects -vH}          & How large the repository really is       \\
\end{keyslong}

\section{Undo and Recovery}\label{app:git-undo}

\begin{keyslong}[0.38]
	\cmd{git restore file}                  & Discard unstaged edits to a file      \\
	\cmd{git restore -{}-staged file}       & Unstage, keeping the edits            \\
	\cmd{git restore -{}-source=HEAD\textasciitilde{}2 f} & Take a file from an old commit \\
	\cmd{git reset -{}-soft HEAD\textasciitilde{}}  & Undo the commit, keep it staged \\
	\cmd{git reset HEAD\textasciitilde{}}           & Undo the commit and unstage    \\
	\cmd{git reset -{}-hard HEAD\textasciitilde{}}  & \textbf{Destroys} the changes  \\
	\cmd{git revert <sha>}                  & A new commit undoing an old one       \\
	\cmd{git clean -n}, \cmd{git clean -fd} & List, then delete untracked files     \\
	\cmd{git reflog}                        & Every position \texttt{HEAD} has held \\
	\cmd{git reset -{}-hard HEAD@\{2\}}     & Go back to before the mistake         \\
	\cmd{git switch -c rescue <sha>}        & Rescue a commit no branch points at   \\
	\cmd{git fsck -{}-lost-found}           & Unreachable objects the reflog missed \\
	\cmd{git filter-repo -{}-path f -{}-invert-paths} & Remove a file from all history \\
\end{keyslong}

\section{Diagnosing Git}\label{app:git-triage}

\begin{keyslong}[0.38]
	``Detached \texttt{HEAD}''             & Harmless. \cmd{git switch -c name} to keep the work \\
	``Your branch is ahead by 2''          & You have unpushed commits             \\
	``no upstream branch''                 & \cmd{git push -u origin <branch>}     \\
	Push rejected, non-fast-forward        & Fetch and rebase; then \cmd{-{}-force-with-lease} if you rewrote \\
	Conflict markers in a file             & Edit, \cmd{git add}, \cmd{git commit}; or \cmd{-{}-abort} \\
	``I reset and lost everything''        & \cmd{git reflog}                       \\
	Uncommitted work lost to \cmd{-{}-hard} & Not recoverable. The one real hazard  \\
	\file{.gitignore} ignored              & The file is already tracked: \cmd{git rm -{}-cached} \\
	Every commit duplicated after a pull   & Someone rewrote a shared branch (\autoref{sec:git-rewrite}) \\
	A secret was pushed                    & Rotate it first; rewrite history second \\
\end{keyslong}

\chapter{Configuration Files}\label{app:config}

\section{A Minimal \texttt{init.lua}}\label{app:config-init}

A self-contained starting point, for when a distribution is more than you want.
Put it in \file{\textasciitilde/.config/nvim/init.lua}; it bootstraps its own
plugin manager on first launch.

\begin{luacode}
-- ── Leaders must be set before any plugin loads ──────────────────────────
vim.g.mapleader = " "
vim.g.maplocalleader = "\\"

-- ── Options ──────────────────────────────────────────────────────────────
local opt = vim.opt
opt.number = true
opt.relativenumber = true
opt.expandtab = true
opt.shiftwidth = 2
opt.tabstop = 2
opt.smartindent = true
opt.ignorecase = true
opt.smartcase = true
opt.undofile = true          -- persistent undo tree
opt.scrolloff = 8
opt.signcolumn = "yes"       -- stop the text jumping when a sign appears
opt.splitright = true
opt.splitbelow = true
opt.termguicolors = true

\end{luacode}

Mappings and autocommands next. \texttt{desc} is what \texttt{which-key} and
\cmd{:map} will show you later.

\begin{luacode}
-- ── Mappings ─────────────────────────────────────────────────────────────
local map = vim.keymap.set
map("n", "<leader>w", "<cmd>write<cr>", { desc = "Save" })
map("n", "<leader>h", "<cmd>nohlsearch<cr>", { desc = "Clear highlight" })
map("n", "<C-h>", "<C-w>h", { desc = "Left window" })
map("n", "<C-j>", "<C-w>j", { desc = "Lower window" })
map("n", "<C-k>", "<C-w>k", { desc = "Upper window" })
map("n", "<C-l>", "<C-w>l", { desc = "Right window" })
map("v", "<", "<gv", { desc = "Dedent, keep selection" })
map("v", ">", ">gv", { desc = "Indent, keep selection" })
map("t", "<Esc><Esc>", "<C-\\><C-n>", { desc = "Leave terminal mode" })

-- ── Autocommands ─────────────────────────────────────────────────────────
vim.api.nvim_create_autocmd("TextYankPost", {
  group = vim.api.nvim_create_augroup("yank_highlight", { clear = true }),
  callback = function() vim.highlight.on_yank({ timeout = 200 }) end,
})

\end{luacode}

Finally the plugin manager, which installs itself the first time Neovim starts.

\begin{luacode}
-- ── Plugins (lazy.nvim, bootstrapped on first run) ───────────────────────
local lazypath = vim.fn.stdpath("data") .. "/lazy/lazy.nvim"
if not (vim.uv or vim.loop).fs_stat(lazypath) then
  vim.fn.system({ "git", "clone", "--filter=blob:none", "--branch=stable",
    "https://github.com/folke/lazy.nvim.git", lazypath })
end
vim.opt.rtp:prepend(lazypath)

require("lazy").setup({
  { "folke/tokyonight.nvim", priority = 1000,
    config = function() vim.cmd.colorscheme("tokyonight") end },
  { "nvim-treesitter/nvim-treesitter", build = ":TSUpdate",
    opts = { highlight = { enable = true }, indent = { enable = true } },
    config = function(_, opts)
      require("nvim-treesitter.configs").setup(opts)
    end },
  { "neovim/nvim-lspconfig" },
  { "lervag/vimtex", ft = "tex", init = function()
      vim.g.vimtex_view_method = "zathura"
    end },
})
\end{luacode}

\section{A Reusable Makefile Skeleton}\label{app:config-makefile}

\section{A \file{.gdbinit} Worth Having}\label{app:config-gdbinit}

\begin{gdbcode}
# ~/.gdbinit
set confirm off                  # stop asking "are you sure"
set pagination off               # no "---Type <return> to continue---"
set history save on              # remember commands between sessions
set history size 10000
set print pretty on              # structs one member per line
set print array on
set print elements 0             # never truncate an array or string
set print null-stop on           # stop a char array at the first NUL
set backtrace limit 100          # tame runaway recursion
set breakpoint pending on        # do not ask about unloaded libraries
set disassembly-flavor intel

# Walk a linked list: plist head
define plist
  set $n = $arg0
  while $n != 0
    printf "%s\n", $n->name
    set $n = $n->next
  end
end
document plist
Walk a linked list from the given head, printing each name.
end

# Trace every call to a function without recompiling: trace parse_line
define trace
  break $arg0
  commands
    silent
    bt 1
    continue
  end
end
document trace
Set a breakpoint that prints one stack frame and carries on.
end
\end{gdbcode}

\section{A \file{.gitconfig} Worth Having}\label{app:config-gitconfig}

\begin{conf}
# ~/.gitconfig
[user]
    name  = Your Name
    email = you@example.com

[init]
    defaultBranch = main

[core]
    editor = nvim
    autocrlf = input          # never store CRLF in the repository

[pull]
    rebase = true             # no "Merge branch 'main' of ..." commits
[rebase]
    autosquash = true         # arrange --fixup commits automatically
    autostash  = true         # stash and restore around a rebase
[push]
    autoSetupRemote = true    # plain `git push' works on a new branch
    followTags = true         # push annotated tags with their commits

[commit]
    verbose = true            # show the diff while writing the message

[diff]
    algorithm = histogram     # better diffs on reordered code
    colorMoved = zebra        # distinguish moved lines from changed ones
[merge]
    conflictStyle = zdiff3    # show the common ancestor in conflict markers

[alias]
    s    = status -sb
    l    = log --oneline --graph --all --decorate
    ll   = log --stat
    d    = diff
    ds   = diff --staged
    co   = switch
    cob  = switch -c
    amend = commit --amend --no-edit
    unstage = restore --staged
    last = log -1 HEAD --stat

# use a different identity for work repositories
[includeIf "gitdir:~/work/"]
    path = ~/.gitconfig-work
\end{conf}

\begin{note}
	\cmd{merge.conflictStyle = zdiff3} is the least known and most useful of these:
	conflict markers then include the \emph{common ancestor} as well as both sides,
	so you can see what each side actually changed rather than guessing from two
	final states.
\end{note}

\chapter{This Machine}\label{app:machine}

Every command in this note was run on one ordinary Linux desktop. Versions
matter more in some places than others --- \cmd{grep} has behaved the same way
for thirty years, while Neovim's LSP API moved twice in two releases --- so they
are recorded here rather than repeated in the text.

\section{Versions}\label{app:machine-versions}

\begin{keyslong}[0.26]
	Operating system & Linux Mint 22.1, kernel 6.8.0                              \\
	\texttt{nvim}    & 0.9.5 --- note \autoref{sec:nvim-lsp}: 0.11 changed the LSP configuration API \\
	\texttt{zsh}     & 5.9 (the login shell)                                       \\
	\texttt{bash}    & 5.2.21                                                      \\
	GNU coreutils    & 9.4                                                         \\
	\texttt{sed}     & GNU sed 4.9 --- \cmd{-i} without an argument is GNU-only    \\
	\texttt{awk}     & GNU Awk 5.2.1                                               \\
	\texttt{make}    & GNU Make 4.3                                                \\
	\texttt{gcc}     & 13.3.0                                                      \\
	\texttt{gdb}     & 15.0.50                                                     \\
	\texttt{git}     & 2.43.0 --- \cmd{git switch} and \cmd{git restore} need 2.23 \\
	\texttt{gh}      & 2.93.0                                                      \\
	Xe\TeX{}         & 3.141592653, TeX Live 2023                                  \\
\end{keyslong}

\section{Installed and Not Installed}\label{app:machine-tools}

Several sections recommend a tool that is not present here. They are named so
that you know what is a recommendation and what was actually exercised.

\begin{keyslong}[0.26]
	Present   & \texttt{rg}, \texttt{jq}, \texttt{tree}, \texttt{htop}, \texttt{tmux}, \texttt{curl}, \texttt{wget}, \texttt{rsync}, \texttt{xclip}, \texttt{zathura}, \texttt{nvr}, \texttt{latexmk} \\
	Present   & GCC's sanitisers --- \cmd{-fsanitize=address,undefined,thread}     \\
	Absent    & \texttt{fd} (\autoref{sec:sh-find}), \texttt{shellcheck} (\autoref{sec:sh-debug}), \texttt{tldr} \\
	Absent    & \texttt{valgrind} and \texttt{rr} (\autoref{sec:gdb-companions})   \\
	Absent    & \texttt{nvim-dap} (\autoref{sec:gdb-frontends}); \texttt{termdebug} ships with Neovim and is used instead \\
	Note      & \cmd{grep} here is \texttt{ugrep} in GNU-compatible mode, not GNU grep \\
	Note      & Core dumps are disabled by default --- \cmd{ulimit -c} reports 0 (\autoref{sec:gdb-core}) \\
\end{keyslong}

\section{The Editor Configuration}\label{app:machine-nvim}

\autoref{ch:vim-setup} describes this in full. In summary: LazyVim on
\texttt{lazy.nvim}, with \texttt{vimtex} (Zathura for viewing, \texttt{nvr} for
inverse search, \texttt{latexmk} with \cmd{-xelatex} into a \file{build/}
directory) and \texttt{LuaSnip} loading snippets from
\file{\textasciitilde/.config/nvim/lua/snippets}. The leader is \key{<Space>}
and the local leader is \key{\textbackslash}, so VimTeX's prefix is
\key{\textbackslash} \key{l}.

\begin{note}
	This note is itself built by the workflow it describes: Xe\TeX{} because the
	template loads \texttt{fontspec}, driven by \cmd{latexmk} from inside Neovim,
	tracked in Git, and typeset from a Makefile of the shape given in
	\autoref{sec:mk-latex}.
\end{note}
